\chapter{La DSL}
\label{cap:dsl}

\begin{intuicion}
La DSL de Graph Hunter es deliberadamente chica: un analista debe
poder escribir una hipótesis en menos de una línea sin abrir
manual. La gramática base es la cadena con flechas; la única
extensión sintáctica reciente (L1 stage 3c) es el separador
\texttt{;} que abre las puertas a \emph{shapes} no-cadena.
\end{intuicion}

\section{Sintaxis chain (legacy)}

\begin{lstlisting}
User -[Auth]-> Host -[Connect]-> IP {k=2} HAVING count(distinct last) >= 5 WITHIN 30m
\end{lstlisting}

\begin{itemize}
    \item \texttt{Type binding}: cada nodo es un \texttt{EntityType}
    seguido opcionalmente por un identificador (\emph{binding}). Ej.\
    \texttt{User u} declara que el primer vértice es de tipo
    \texttt{User} y se referenciará como \texttt{u} en clauses
    posteriores.
    \item \texttt{-[RelationType\{predicates?\}]->}: arista con
    \emph{relation type} y predicados opcionales sobre la metadata.
    \item \texttt{\{key=value\}} sobre el \emph{dest type}: predicados
    sobre el nodo destino.
    \item \texttt{\{k=N\}} al final: \emph{k-simplicity} (cuántas veces
    puede repetirse un vértice en el \emph{path}).
    \item \texttt{HAVING count(distinct col) >= N}: agregación
    post-\emph{hoc}.
    \item \texttt{WITHIN Ns/Nm/Nh}: ventana temporal entre \emph{first}
    y \emph{last edge} del \emph{path}.
\end{itemize}

\section{Sintaxis Graph (L1 stage 3c)}

\begin{lstlisting}
User c -[Auth]-> Host h1; c -[Auth]-> Host h2; c -[Auth]-> Host h3
\end{lstlisting}

\begin{intuicion}
El primer \texttt{;} marca la hipótesis como \emph{shape}
\texttt{Graph}: cada \texttt{(step, side)} es su propia variable a
menos que un \emph{binding} explícito una posiciones. Una cláusula
después de \texttt{;} puede arrancar de dos formas:

\begin{enumerate}
    \item \emph{Bare binding reference}: un identificador que ya fue
    declarado en una cláusula anterior (\texttt{c} arriba). El motor
    hereda el \texttt{EntityType} del binding original.
    \item \emph{Fresh \texttt{Type binding} declaration}: introduce
    un vértice independiente.
\end{enumerate}
\end{intuicion}

\subsection{Disambiguación}

\begin{itemize}
    \item Identificadores que coinciden con un \texttt{EntityType}
    builtin (\texttt{User}, \texttt{Host}, \texttt{IP}, etc.) siempre
    se parsean como \emph{type declaration}, sombreando bindings del
    mismo nombre.
    \item Identificadores desconocidos producen un error con un
    \emph{hint} listando los bindings declarados.
\end{itemize}

\subsection{Estructura interna del parser}

\begin{figure}[H]
\centering
\begin{tikzpicture}[
    node distance=0.5cm,
    every node/.style={font=\scriptsize, align=center},
    state/.style={draw=accent, thick, rounded corners=1mm,
                  minimum height=0.7cm, minimum width=2.6cm,
                  fill=accentlight}
]
\node[state] (start) {\texttt{parse()}};
\node[state, below=of start] (head) {\texttt{parse\_entity\_type\_with\_binding}};
\node[state, below=of head] (chain) {\texttt{parse\_chain\_steps}};
\node[state, below=of chain] (sep) {peek \texttt{;}};
\node[state, below left=of sep] (clause) {\texttt{parse\_clause\_head}};
\node[state, below right=of sep] (trail) {\texttt{parse\_trailing\_clauses}};

\draw[->, >=Latex, accent, thick] (start) -- (head);
\draw[->, >=Latex, accent, thick] (head) -- (chain);
\draw[->, >=Latex, accent, thick] (chain) -- (sep);
\draw[->, >=Latex, accent, thick] (sep) -- node[left]{si \texttt{;}} (clause);
\draw[->, >=Latex, accent, thick, dashed]
    (clause.east) to[out=0,in=180] (chain.west);
\draw[->, >=Latex, accent, thick] (sep) -- node[right]{else} (trail);
\end{tikzpicture}
\caption{Flujo del parser. \texttt{parse\_chain\_steps} es el bucle
interior reusable; cada \texttt{;} dispara un re-entry vía
\texttt{parse\_clause\_head}, que decide entre \emph{bare binding ref}
y \emph{type declaration}.}
\label{fig:dsl-parser-flow}
\end{figure}

Cómo se lee el flujo en cada nodo del diagrama:

\begin{itemize}
    \item \texttt{parse\_entity\_type\_with\_binding} (\emph{first
    clause head}): consume un \texttt{Type} obligatorio y, si lo
    sigue un identificador no reservado, lo registra como
    \emph{binding}. Es la única forma de declarar el primer vértice
    --- las hipótesis no pueden empezar con un \emph{bare binding ref}
    porque no hay nada declarado todavía.
    \item \texttt{parse\_chain\_steps}: bucle que consume \emph{N}
    pasos \texttt{-[rel]-> Type binding} encadenados. Termina cuando
    ve EOF, \texttt{\{} (suffix \texttt{\{k=N\}} o predicate block),
    \texttt{;} (clause separator), o un identificador
    alfabético (HAVING/WITHIN). Cada terminación delega al manejador
    correspondiente.
    \item \texttt{parse\_clause\_head} (\emph{tras un \texttt{;}}):
    peeka el identificador líder y decide:
    \begin{itemize}
        \item Si matchea un \emph{builtin type name}
        (\texttt{User}, \texttt{Host}, \texttt{IP}, etc.) $\Rightarrow$
        rewind y delega a \texttt{parse\_entity\_type\_with\_binding}
        --- introduce un vértice fresco con su propio binding.
        \item Si matchea un \emph{binding} ya declarado en cláusulas
        anteriores $\Rightarrow$ retorna
        \texttt{(existing\_type, Some(name))} para que el bucle de
        \texttt{parse\_chain\_steps} arranque desde ese vértice
        compartido.
        \item Si no es ninguno de los dos $\Rightarrow$ error con
        \emph{hint} listando los \emph{bindings} válidos. La opción
        de fallback a \texttt{Other(\dots)} se descarta a propósito
        para no swallow-tear errores de tipeo del analista.
    \end{itemize}
    \item \texttt{parse\_trailing\_clauses}: consume \emph{HAVING} y
    \emph{WITHIN} en cualquier orden hasta encontrar \texttt{\{k=N\}}
    o EOF. Estos no pueden re-entrar el bucle de chain steps; son
    estrictamente \emph{post-fix} a la última cláusula del cuerpo.
\end{itemize}

\section{Del AST al matcher}

El AST que el parser produce es un \texttt{Hypothesis}:

\begin{lstlisting}[language=Rust]
pub struct Hypothesis {
    pub name: String,
    pub steps: Vec<HypothesisStep>,
    pub k_simplicity: usize,
    pub aggregations: Vec<AggClause>,
    pub window_seconds: Option<i64>,
    pub shape: HypothesisShape,   // Chain | Graph
}
\end{lstlisting}

El \texttt{shape} es la pieza nueva que destrabó L1 stage 3.
El \emph{matcher} branchea en \texttt{search\_temporal\_pattern\_yannakakis}:

\begin{itemize}
    \item \texttt{Chain}: \texttt{semijoin\_reduce\_chain}
    \(+\) \texttt{enumerate\_chain}.
    \item \texttt{Graph}: \texttt{JoinTree::build}
    \(+\) \texttt{enumerate\_tree}.
\end{itemize}

\section{Plug-and-play en todos los \emph{transports}}

La DSL está expuesta vía \texttt{parse\_dsl} en
\texttt{platform/dsl}. Todos los \emph{transports} la consumen sin
flag adicional:

\begin{itemize}
    \item HTTP: \texttt{POST /run\_hunt} con \texttt{\{ "dsl": "..." \}}.
    \item MCP: \texttt{cmd\_parse\_dsl} +
    \texttt{cmd\_run\_hunt}.
    \item Tauri: \texttt{cmd\_run\_hunt(hypothesis\_json: String)}
    (el frontend serializa el AST).
\end{itemize}

Un cliente que envía una hipótesis con \texttt{;} hoy mismo recibe
una \texttt{Hypothesis::shape = Graph} que el \emph{planner} dispatcha
al \emph{tree-walk} sin tocar nada más.
